% ============================================================
%  Введение
% ============================================================
\chapter*{Введение}
\addcontentsline{toc}{chapter}{Введение}

\section*{Актуальность работы}

Экзотические деривативы~--- инструменты, выплата по которым зависит не
только от спотовой цены базового актива в момент исполнения, но и от
всей траектории цен на отрезке жизни контракта~--- занимают значительную
долю внебиржевого рынка производных финансовых инструментов. Они активно
используются для тонкой настройки риск-профиля портфеля: хеджирования
средних цен (азиатские опционы), защиты от резких движений
(барьерные опционы), управления риском по корзинам активов (радужные
опционы) и торговли структурой корреляций (корреляционные свопы).

При этом для большинства таких инструментов отсутствуют удобные
аналитические формулы цены и греков. Стандартным методом оценки
становится метод Монте-Карло, который, однако, плохо масштабируется при
большом числе одновременно котируемых контрактов и при необходимости
часто пересчитывать греки для динамического хеджирования. Это создаёт
естественный запрос на использование методов машинного обучения, которые
позволяют не только аппроксимировать цену, но и сразу получать частные
производные за счёт автоматического дифференцирования~\cite{joas2024}.

\section*{Цель работы}

Цель настоящей работы~--- разработать и сравнить методы оценки опционов,
выплата по которым зависит от траектории, объединяющие классические
финансово-математические подходы (модели стохастической волатильности и
корреляции, декомпозиция цены, хеджирующие портфели) с методами
машинного обучения для оценки оста\-точного риска.

\section*{Задачи}

\begin{enumerate}
  \item Систематизировать классы экзотических деривативов и обозначить
        границы применимости известных аналитических методов.
  \item Изложить теорию корреляционных свопов, в том числе модели
        возврата к среднему, Якоби и Фишера--OU, и обсудить их
        применение к хеджированию рынка от системных пузырей.
  \item Сформулировать декомпозицию цены rainbow-опциона на
        аналитически считаемые компоненты и нелинейный остаток.
  \item Построить статическую и динамическую модели хеджирования
        rainbow-опциона; оценить роль транзакционных издержек при
        ребалансировке.
  \item Обучить модели машинного обучения для оценки нелинейной части
        rainbow-опциона на основе сгенерированного датасета и сравнить
        с прямой оценкой и с декомпозиционным подходом.
  \item Сравнить методы по метрикам RMSE и MAE.
\end{enumerate}

\section*{Структура работы}

Работа состоит из введения, пяти глав, заключения, списка литературы
и приложений. Глава~\ref{ch:exotic} посвящена обзору экзотических
деривативов и методам их оценки. Глава~\ref{ch:corr} подробно
рассматривает корреляционные свопы и rainbow-опционы. В
главе~\ref{ch:ml} обсуждается применение методов машинного обучения.
Глава~\ref{ch:data} описывает использованные данные. В
главе~\ref{ch:results} приводятся результаты вычислительных
экспериментов и их сравнение.

\newpage
